%% ─────────────────────────────────────────────────────────────────────────────
%% CLASE emate-ucr  —  prueba corta
%% ─────────────────────────────────────────────────────────────────────────────
%%
%% Este archivo muestra los comandos adicionales disponibles para
%% evaluaciones formales (pruebas cortas, exámenes).
%% Todos son opcionales: si no se usan, el documento queda igual
%% que una hoja de ejercicios sencilla.
%%
\documentclass{emate-ucr}
% \documentclass[soluciones]{emate-ucr}   % <-- versión con soluciones

%% ── Metadatos obligatorios ───────────────────────────────────────────────────
\curso{I Ciclo 2026\\MA-1022\\Cálculo para Ciencias Económicas II} % encabezado de página
\encabezado{Prueba Corta 1}              % título del documento

%% ── Metadatos opcionales de evaluación ───────────────────────────────────────
%%
%% Los tres comandos siguientes son opcionales e independientes entre sí.
%% Si se definen, \imprimirtitulo los muestra en una línea adicional
%% centrada bajo el título, separada por una regla:
%%
%%   Fecha: 20 de marzo de 2026  |  Duración: 50 minutos  |  Valor: 15 puntos
%%
%% Si ninguno está definido, esa línea no aparece (retrocompatible con
%% hojas de ejercicios que no los usan).
%%
\fecha{20 de marzo de 2026}       % fecha de aplicación
\duracion{50 minutos}             % tiempo disponible
\valor{\totalpuntos\ puntos}      % valor total (compilar dos veces)
\nombreejercicio{Problema}        % palabra usada en lugar de "Ejercicio"

\begin{document}

%% \imprimirtitulo
%%   Imprime el título y, si corresponde, la línea de metadatos.
%%   Llamar siempre al inicio del documento.
%%
\imprimirtitulo

%% \datosestudiante
%%   Imprime una línea con espacios para que el estudiante escriba
%%   su nombre y número de carné:
%%
%%     Nombre: _______________________________  Carné: __________
%%
%%   Es opcional: omitirlo en hojas de ejercicios que no requieren
%%   identificación individual.
%%
\datosestudiante

%% ── Entorno instrucciones ─────────────────────────────────────────────────────
%%
%% Imprime una caja enmarcada con las indicaciones de la evaluación.
%% Cada indicación va con \item (lista numerada automáticamente).
%% Usar en lugar del entorno indicaciones cuando se quiere un formato
%% más formal con borde visible.
%%
%% (También existe \begin{indicaciones}...\end{indicaciones} para un
%%  formato más simple, sin caja: solo la palabra "Indicaciones:" en
%%  negrita seguida del texto libre, sin lista numerada.)
%%
\begin{instrucciones}
  \item Esta prueba es individual y con libro cerrado.
  \item Justifique cada respuesta. Una respuesta sin justificación no se
        calificará.
  \item Si necesita más espacio, use el reverso de la hoja e indíquelo
        claramente.
  \item Puede usar calculadora científica (no graficadora).
\end{instrucciones}

\begin{ejercicio}[10]
  Resuelva el siguiente sistema de ecuaciones aplicando eliminación de
  Gauss-Jordan. Dé el conjunto solución.
  \[
    \begin{cases}
       x + 2y +  z =  5 \\
      2x + 5y + 2z = 11 \\
           y +  z =  3
    \end{cases}
  \]

  \begin{solucion}
    La matriz aumentada es
    \[
      \left(\begin{array}{ccc|r} 1 & 2 & 1 & 5 \\ 2 & 5 & 2 & 11 \\ 0 & 1 & 1 & 3
      \end{array}\right)
      \xrightarrow{f_2\to f_2-2f_1}
      \left(\begin{array}{ccc|r} 1 & 2 & 1 & 5 \\ 0 & 1 & 0 & 1 \\ 0 & 1 & 1 & 3
      \end{array}\right)
      \xrightarrow{f_3\to f_3-f_2}
      \left(\begin{array}{ccc|r} 1 & 2 & 1 & 5 \\ 0 & 1 & 0 & 1 \\ 0 & 0 & 1 & 2
      \end{array}\right).
    \]
    De la fila~3: $z=2$. De la fila~2: $y=1$. De la fila~1: $x=1$.
    El conjunto solución es $S=\{(1,\,1,\,2)\}$.
  \end{solucion}
\end{ejercicio}

\begin{ejercicio}[5]
  Determine el rango de la matriz
  $\begin{pmatrix} 1 & 2 & 3 \\ 2 & 5 & 7 \\ 1 & 3 & 4 \end{pmatrix}$.

  \begin{solucion}
    Reduciendo a forma escalonada se obtienen 2 filas no nulas.
    El rango es $\mathbf{2}$.
  \end{solucion}
\end{ejercicio}

\end{document}
